\documentclass{llncs}

%\usepackage{llncsdoc}

%\usepackage{makeidx}  % allows for indexgeneration
\usepackage{graphicx}
\usepackage[T1]{fontenc}
\usepackage[english]{babel}
\usepackage[utf8]{inputenc}
\usepackage{multirow}

\usepackage{url}
\usepackage{rotating}

%%%Math
\usepackage{latexsym}
% \usepackage{amsmath}
% \usepackage{amssymb}
% \usepackage{amsthm}
%\usepackage{eurosans}

\usepackage{eurosym}

\usepackage{longtable}

\usepackage{listings}

\usepackage{color}
\usepackage{textcomp}


\definecolor{gray}{gray}{0.5}
\definecolor{green}{rgb}{0,0.5,0}


\begin{document}
\title{A natural language and semantic-based technique to unify corporate names in the e-Procurement sector. The CORFU approach.\thanks{THANKS}}

\titlerunning{}

\author{Jose Mar\'{i}a Alvarez-Rodr\'{i}guez\inst{1}} 

\authorrunning{Jose Mar\'{i}a Alvarez-Rodr\'{i}guez}


\tocauthor{Jose Mar\'{i}a Alvarez}


\institute{The South East European Research Center\\   
  \email{{jmalvarez@seerc.org}},\\
   WWW home page: \texttt{http://www.seerc.org}, \\
}


\date{}

\maketitle

\renewcommand{\labelitemi}{$\bullet$}

\begin{abstract}
Public administrations are currently facing a big challenge trying to improve both the peformance and the transparency of administrative processes.
In this context the e-Government and Open Linked Data initiatives are tackling existing interoperability and 
integration issues among ICT-based systems but the creation of a real transparent environment requires 
much more than the simple publication of data and information in specific open formats; data and information 
quality is the next major step in the pubic sector. More specifically in the e-Procurement domain there is a 
vast amount of valuable data that is already available via the Internet protocols and formats and can be used 
for the creation of new added-value services. Neverthless the simple extraction of statistics or creation of reports 
can imply extra tasks with regards to clean, prepare and reconcile data. 
From a transparency point of view one of the most interesting services lies in tracking rewarded contracts (type, location, and supplier). 
Depending on the capabilities of the public organization this kind of basic report can turn into a 
complex task due to a lack of standardization in supplier names or the use of different descriptors for the type of contract. That is why 
this paper presents a stepwise method based on natural language processing and semantics to deal with the hetereogenities in corporate names. 
Furthermore a research study to evaluate the precision and recall of the proposed technique, using as use case the public dataset of rewarded public 
contracts in Australia during the period 2004-2012, is also presented. Finally some discussion, 
conclusions and future work are also outlined.
\end{abstract}

\section{Introduction}
Public bodies are continuously publishing procurement opportunities in which 
valuable metadata is available. Depending on the stage of the process new data arises such 
as the supplier name that has been rewarded with the public contract. In this 
context the extraction of statistics on how many contracts have been 
rewarded to the same company is a relevant indicator to evaluate the transparency 
of the whole process. Although companies that want to tender for a public contract must be 
oficially registered and have an unique identification number, the truth is 
that in most of rewarded contracts the supplier is only identified by a name or a string literal typed 
by a civil-servant. In this sense there is not usually a connection between 
the oficial company registry and the process of rewarding contracts implying different 
naming problems and data inconsistencies that are spread to next stages preventing future 
activities such as reporting.

In the case of the type of contract and location, there are already standardized~\cite{DBLP:journals/ijseke/AlvarezLSASL12} product 
scheme classifications such as the Common Procurement Vocabulary (2003 and 2008), the Combined Nomenclature (2012), 
the Central Product Classificationby in the European Union, the International Standard Industrial Classification of 
All Economic Activities (Rev. 4) in the United Nations or the North American Industry Classification System (2007 and 2012) 
in the Government of United States are currently used with different objectives such as statistics, tagging or 
information retrieval. Geolocated information can be also found in different common datasets and nomenclatures such as 
the Nomenclature of territorial units for statistics in the European Union, the Geonames dataset~\footnote{\url{http://www.geonames.org/}}, the GeoLinkedData 
initiative~\cite{DBLP:conf/dexa/Lopez-PellicerSCZM10} or the traditional list of countries and ISO-codes.

However corporate, organization, firm, company or institution names (hereafter these names will be used to refer to 
the same entity) and structure is not yet standardized at global scope and only some classifications of economic activities or 
company identifiers can be found such as the TARIC (Online customs tariff database ) database. Thus the simple task of grouping contracts by a supplier is not a mere process of 
searching by the same literal. Technical issues such as hyphenation, use of abbreviations or acronynms an transliteration are common problems that must be addressed in order 
to provide a final corporate name. Existing works in the field of Name Entity Recognition~\cite{citeulike:1657521} (NER) or 
name entity disambiguation~\cite{Sarmento:2009:AWN:1602022.1602085,Klein:2003:NER:1119176.1119204} have already addressed these issues. 
Neverthless the problem that is being tackled in these approaches lies in the identification of organization names in 
a raw text while in the e-Procurement sector the string literal identifying a supplier is already known.

In the particular case of the Australian e-Procurement domain, the supplier name is introduced by typing a string literal without any assistance or 
auto-complete method. Obviously a variety of errors and variants for the same company, see Table~\ref{tabla:aus-suppliers}, 
can be found such as misspelling errors~\cite{NorvigSpelling,StanfordSpelling}, name and acronym mismatches~\cite{Yeates99automaticextraction,Ratinov:2004:AES:1025132.1026366} 
or context-aware data that is already known when the dataset is processed, e.g. country or year. Furthermore it is also well-known 
that a large company can be divided into several divisions or departments but from a statistical point of view grouping data by a supplier name 
should take into account all rewarded contracts regardless the structure of the company.

On the other hand the application of semantic technologies and the Linking Open Data initiative (hereafter LOD)~\cite{Berners-Lee-2006,Heath_Bizer_2011}  
in several fields like e-Government (e.g. the Open Government Data effort) tries to improve the knowledge about a specific area providing 
common data models and formats to share information and data between agents. More specifically, in the European e-Procurement 
context~\cite{e-Proc-map-paper} there is an increasing commitment to boost the use of electronic communications and transactions 
processing by government institutions and other public sector organizations in order to provide added-value services with special focus on SMEs. 
More specifically the LOD initiative seeks for creating a public data realm in which one the principles of this initiative that lies in the 
unique identification or resources through URIs can become real. Thus entity reconciliation techniques~\cite{Serimi,conf/www/MaaliCP11} 
coming from the ontology mapping and alignment areas or algorithms based on Natural Language Processing (hereafter NLP) have been 
designed to link similar resources already available in different vocabularies, datasets or databases such as DBPedia or Freebase. 
Nevertheless the issue of unifying supplier names as a human would do faces new problems that have been tackled in 
other research works~\cite{Galvez2006} to extract statistics of performance in bibliographic databases. The main objective is not just a 
mere reconciliation process to link to existing resources but to create a unique name or link ($n$ string literals $\to$ $1$ company $\to$ $1$ URI). 
For instance in the case of the ongoing example the string literals ``Oracle'' and ``Oracle University'' could be respectively aligned to the entity $<$Oracle\_Corporation$>$ and $<$Oracle\_University$>$ but 
the problem of grouping by a unique (\textit{Big}) name, identifier or resource still remains. That is why a context-aware method based on NLP 
techniques combined with semantics has been designed, customized and implemented trying to exploit the naming convention of a specific dataset with the aim 
of grouping $n$ string literals $\to$ $1$ company and, thus, easing the next natural process of entity reconciliation.

The remainder of this paper is structured as follows. Section 2 a literature review is presented. Afterwards next section outlines main mismatches in corporate names. Section 4 presents 
the CORFU approach to unify corporate names. The evaluation section exposes and discusses the experimentation carried out to test the presented approach using as a dataset the rewarded 
contracts of Australia in the period 2004-2012. Finally conclusions summarizes the main outcomes of this work 
and some open issues are also presented as future work.


\begin{table}[!htb]
\renewcommand{\arraystretch}{1.3}
\begin{center}
\begin{tabular}{|p{7cm}|p{7cm}|}
\hline
  \textbf{Raw Supplier Name} & \textbf{Target (potential) Supplier Name and URI}  \\  \hline
  ``Accenture'' & \multirow{12}{*}{``Accenture''} \\
  ``ACCENTURE'' & \multirow{12}{*}{\scriptsize\url{http://live.dbpedia.org/resource/Accenture}} \\ 
  ``Accenture Aust Holdings'' & \\  
  ``Accenture Aust Holdings Pty Ltd'' & \\
  ``Accenture Australia'' & \\
  ``ACCENTURE AUSTRALIA'' & \\
  ``Accenture  Australia Holding P/L'' & \\
  ``Accenture Australia Holdings P/Ltd'' & \\
  ``Accenture Australia Holdings Pty'' & \\
  ``Accenture Australia Holdings Pty Lt''  & \\  
  ``Accenture Australia Limited'' & \\
  \ldots  & \\
  ``Accenture Australia Ltd'' & \\ \hline
  ``Microsoft Australia'' & \multirow{3}{*}{``Microsoft''} \\
  ``Microsoft Australia Pty Ltd'' & \multirow{3}{*}{\scriptsize\url{http://live.dbpedia.org/resource/Microsoft}} \\
  \ldots  & \\
  ``Microsoft Enterprise Services'' & \\ \hline
  ``Oracle (Corp) Aust Pty Ltd''  & \multirow{14}{*}{``Oracle''} \\
  ``Oracle Corp (Aust) Pty Ltd''  & \multirow{14}{*}{\scriptsize\url{http://live.dbpedia.org/resource/Oracle_Corporation}} \\
  ``Oracle Corp Aust Pty Ltd'' & \\
  ``ORACLE CORP AUST PTY LTD'' & \\
  ``Oracle Corp. Australia'' & \\
  ``ORACLE CORP AUSTRALIA P/L'' & \\
  ``Oracle Corp. Australia Pty.Ltd.'' & \\
  ``ORACLE CORP AUSTRALIA PTY LTD'' & \\
  ``Oracle Corpoartion (Aust) Pty Ltd'' & \\
  ``Oracle Corporate Aust Pty Ltd'' & \\
  ``Oracle Corporation'' & \\
  ``Oracle Risk Consultants'' & \\
  ``ORACLE SYSTEMS (AUSTRALIA) PTY LTD'' & \\
  \ldots  & \\
  ``Oracle University''  & \\ \hline
  ``PRICEWATERHOUSECOOPERS(PWC)''  & \multirow{8}{*}{``PricewaterhouseCoopers''} \\
  ``PricewaterhouseCoopers Securities Ltd''& \multirow{8}{*}{\scriptsize\url{http://dbpedia.org/resource/PricewaterhouseCoopers}} \\
  ``PriceWaterhouseCoopers Securities Ltd'' & \\
  ``PRICEWATERHOUSE COOPERS SECURITIES LTD'' & \\
  ``PricewaterhouseCoopers Services LLP'' & \\
  ``Pricewaterhousecoopers Services Pty Ltd'' & \\
  ``PriceWaterhouseCoopers (T/A: PriceWaterhouseCoopers Legal)'' & \\
  \ldots  & \\
  ``Pricewaterhouse (PWC)'' & \\ \hline
  \ldots & \ldots \\
  \hline
  \end{tabular}
  \caption{Examples of supplier names in Australian rewarded contracts.}
  \label{tabla:aus-suppliers}
  \end{center}
\end{table} 


\section{Related Work}
The literature review of this work covers the following areas:
\begin{itemize}
 \item Natural Language Processing and Computational Linguistics. In these research areas common works dealing with the aforementioned data hetereogenities 
   such as misspelling errors~\cite{NorvigSpelling,StanfordSpelling} and name/acronym mismatches~\cite{Yeates99automaticextraction,Ratinov:2004:AES:1025132.1026366}, 
  in the lexical, syntactic and semantic levels can be found. These approaches can be applied to solve general problems and usually follow a 
  traditional approach of text normalization, lexical analysis, pos-tagging word according to a grammar and semantic analysis to filter or 
  provide some kind of service such as information/knowledge extraction, reporting, sentiment analysis or opinion mining. 
  Well-stablished APIs such as NLTK~\cite{LoperBird02} for Python, Lingpipe~\cite{Lingpipe}, OpenNLP~\cite{OpenNLP} or Gate~\cite{Gate} for Java, WEKA~\cite{read12:_scalab} 
  (a data mining library with NLP capabilities), the Apache Lucene and Solr~\cite{rafa2011apache} search engines provide the proper building blocks to build natural-language based applications. 
  Recent times have seen how the analysis of social networks such as Twitter~\cite{Li:2012:TNE:2348283.2348380,Gimpel:2011:PTT:2002736.2002747}, the extraction of 
  clinical terms~\cite{Wang:2009:ARN:1667884.1667888} for electronic health records, the creation of bibliometrics~\cite{Galvez2006,Morillo:2013:TAA:2424697.2424727} or 
  the idenfication of gene names~\cite{Krauthammer:2004:TIB:1053007.1053018,Galvez2012} to name a few have tackled the problem of entity recognition and extraction from raw sources. 
  Other supervised techniques~\cite{Bohn:2006:PHD} have also be used to train data mining-based algorithms with the aim of creation classifiers.
 
 \item Semantic Web. More specifically in the LOD initiative~\cite{Berners-Lee-2006} the use of entity reconciliation techniques to uniquely identify resources 
 is being currently explored. Thus an entity reconciliation pocess can be briefly defined as the method for looking and mapping~\cite{DBLP:conf/semweb/IseleJB10,DBLP:conf/icwe/IseleJB12} two different 
 concepts or entities under a certain threshold. There are a lot of  works presenting solutions about concept mapping, entity reconciliation, etc. 
 most of them are focused on the previous NLP techniques~\cite{conf/www/MaaliCP11,Serimi} (if two concepts have similar literal descriptions then they should be similar) 
 and others (ontology-based) that also exploit the semantic information (hierarchy, number and type of relations) to establish a potential mapping 
 (if two concepts share similar properties and similar super classes then these concepts should be similar). Apart from that 
 there are also machine learning techniques to deal with these mismatches in descriptions using statistical approaches. Recent times, 
 this process has been widely studied and applied to the field of linking entities in the LOD realm, for instance using the DBPedia~\cite{Mendes:2011:DSS:2063518.2063519}. 
 Although there is no way of automatically creating a mapping with a 100\% of confidence (without human validation) a mapping under a certain percentage of confidence can be 
 enough for most of user-based services such as visualization. However, in case of using these techniques as previous step of a reasoning or 
 a formal verification process this ambiguity can lead to infer incorrect facts and must be avoided without a previous human validation. 

 On the other hand the use of semantics is also being applied to model organizational structures. In this case the notion 
 of corporate is presented in several vocabularies and ontologies as Dave Reynolds (Epimorphics Ltd) reports~\footnote{\url{http://www.epimorphics.com/web/wiki/organization-ontology-survey}}. 
 Currently the main effort is focused in the designed of the Organizations Vocabulary (a W3C Working Draft) in which the structure and 
 relationships of companies are being modelled. This proposal is especially relevant in the next aspects:  
 1) to unify existing models to provide a common specification; 2) to apply semantic web technologies and the Linked Data approach to enrich 
 and publish the relevant corporate information; 3) to provide access to the information via standard protocols 
 and 4)to  offer new services that can exploit this information to trace the evolution and behavior of the organization over time.

 
 \item Corporate Databases. Although corporate information such as identifier, name, economic activity, contact person, address or 
 financial status is usually publicly available in the official government registries the access to this valuable information can be 
 tedious due to different formats, query languages, etc. That is why other companies have emerged trying to index and exploit 
 these public repositories; selling reporting services that contain an aggregated version of the corporate information. Taking as 
 an example the Spanish realm, the Spanish Chambers of Commerce~\footnote{\url{http://www.camerdata.es/php/eng/fichero\_empresas.php}},
 Empresia.es~\footnote{\url{http://empresia.es}} or Axesor.es~\footnote{\url{http://www.axexor.es}} manage a database of companies and individual 
 entrepreneurs. This situation can be also transpose to the international scope, for instance Forbes keeps a list of 
 the most representative companies in different sectors. The underlying problems lies in the lack of unique identification, same company data in more 
 than a source, name standardization, etc. and, as a consequence, difficulty of tracking company activity. In order to tackled these problems some 
 initiatives applying the LOD principles such as the Orgpedia~\footnote{\url{http://tw.rpi.edu/orgpedia/}} in United States or 
 ``The Open Database Of The Corporate World''~\footnote{\url{http://opencorporates.com/}} have scrapped and published the information 
 of companies creating a large database containing ($54,080,317$ of companies in May 2012) with high-valuable information like the company 
 identifier, e.g. Figure~\ref{figure:open} shows an example of an ``OpenCorporates'' organization. 
 Apart from that, reconciliation services have also been provided but the problem of mapping ($n$ string literals $\to$ $1$ company $\to$ $1$ URI, 
 as a human would do and the previous section has presented) still remains. Finally public web sites and major social networks such as Google 
 Places, Google Maps, Foursquare, Linkedin Companies, Facebook or Tuenti provide APIs and information managed by the own companies that is supposed 
 to be specially relevant to enrich existing corporate data once a company is uniquely identified.
 
  
 
\begin{figure}[!h]
\begin{center}
\begin{lstlisting}[language=SPARQL]
...
<http://opencorporates.com/id/companies/us_az/F07503757#id> 
	dct:created "1995-06-01"^^xsd:date;
	a adms:Identifer;
	skos:notation "F07503757";
	adms:schemaAgency "Arizona Corporation Commission".

<http://opencorporates.com/id/companies/us_az/F07503757#ra> 
	a locn:Address;
	locn:fullAddress "% CORPORATION SERVICE COMPANY, 
	2338 W ROYAL PALM RD STE-J, PHOENIX, AZ 85021".

<http://opencorporates.com/id/companies/us_az/F07503757> 
	opencorporates:companyType "CORPORATION";
	opencorporates:legalName "ORACLE SOFTWARE CORPORATION (FN)";
	a <http://s.opencalais.com/1/type/er/Company>,
		legal:LegalEntity;
	rdfs:label "ORACLE SOFTWARE CORPORATION (FN)";
	vCard:adr _:bnode1324364416;
	legal:companyType "CORPORATION";
	legal:legalIdentifier 
	  <http://opencorporates.com/id/companies/us_az/F07503757#id>;
	legal:legalName "ORACLE SOFTWARE CORPORATION (FN)";
	locn:registeredAddress 
	  <http://opencorporates.com/id/companies/us_az/F07503757#ra>.
...
\end{lstlisting}
\caption{Partial Information in the N3 format about an ``Oracle'' company in ``Open Corporates''.}
\label{figure:open}
\end{center}
\end{figure}
 
 
\end{itemize}

 \section{The CORFU approach}
According to~\cite{Galvez2006,Morillo:2013:TAA:2424697.2424727} institutional name variations can be 
classified into two different groups: 1) Non-acceptable variations (affect to the meaning) due to mispelling or translation errors and 
2) acceptable variations (do not affect to the meaning) that correspond to different syntax forms such as abbreviations, use of acronyms or contextual 
information like country, suborganization, etc. In order to address these potential variations the CORFU (Company, ORganization and Firm Unifier) approach 
seeks for providing a stepwise method to unify corporate names using NLP and semantics based techniques as a previous step to perform 
an entity reconciliation process. The execution of CORFU comprises several common but customized steps in natural language processing applications such as 
1) text normalization; 2) filtering; 3) comparison and clusterization and 4) linking to an existing information resource. The CORFU unifier 
make an intensive use of the Python NLTK API and other packages for querying REST services or string comparison. Finally and 
due to the fact that the corporate name can change in each step the initial raw name must be saved as well as contextual information such as dates, acronyms or locations. 
Thus common contextual information can be added to create the final unified name.

\begin{enumerate}
 \item Normalize raw text and remove duplicates. This step is comprised of: 1) remove strange characters and punctuation marks but keeping 
 those that are part of a word avoiding potential changes in  abbreviations or acronyms; 2) lowercase the raw text (although some semantics can 
 be lost previous works and empirical tests show that this is the best approach); 3) remove duplicates and 4)lemmatize the corporate name. 
 The implementation of this step to clean the corporate name has ben performed using a combination of the aforementioned API and 
 the Unix scripting tools AWK and SED. In this case, Figure~\ref{figure:step-1} presents 
 two relevant snippets of code for cleaning the name and making a basic word normalization. 
 
\begin{figure}[!h]
\begin{center}
\begin{lstlisting}[language=Python]        
rawname = filter(lambda x: x in string.letters or 
	      x in string.whitespace, line)
...
def normalize(self, word):
  word = word.lower()
  word = self.lemmatizer.lemmatize(word)
  return word
\end{lstlisting}
\caption{Normalization and data cleansing using the Python NLTK API.}
\label{figure:step-1}
\end{center}
\end{figure}


   
\item Filter the basic set of common stopwords in English. A common practice in NLP relies in the construction 
of stopwords sets that can filter some non-relevant words. Nevertheless the use of this technique must 
consider two keypoints: 1) there is a common set of stopwords for any language than can be often used as a filter and 
2) depending on the context the set of stopwords should change to avoid filtering relevant words. In this particular 
case, a common and minimal set of stopwords in English provided by NLTK has been used. Thus the normalized corporate name is 
transformed into a new set of words. Figure~\ref{figure:step-2} presents the function for removing a set of 
words given a another set, it can also be applied to other stages that require filtering capabilities.

\begin{figure}[!h]
\begin{center}
\begin{lstlisting}[language=Python] 
from nltk.corpus import stopwords
self.stop_words_wn = Set(stopwords.words('english'))
...
def remove_set(self, set,  name): 
  token_names= word_tokenize(name)       
  filtered_token_list = [w for w in  token_names if not w in set ]
  cleaned_name = " ".join(["".join(filtered_token) 
    for filtered_token in filtered_token_list])
  return cleaned_name
  
stop_unified_name = self.remove_set(self.stop_words_wn, name)
\end{lstlisting}
\caption{Filtering words with the Python NLTK API.}
\label{figure:step-2}
\end{center}
\end{figure}

\item Filter the expanded set of most common words in the dataset. Taking into account 
the aforementioned step this stage is based on the construction of a customized stopwords 
set for corporate names that is also expanded with Wordnet (ver. 3.0) synonyms with the aim of 
exploiting semantic relationships. In order to create this set two strategies have been followed and applied: 1) create the set 
of words by hand (accurate but very time-consuming) and 2) extract automatically the set of ``most common words'' from the 
working dataset and make a hand-validation (less accurate and time-consuming). Figure~\ref{figure:step-3} partially 
shows these approaches implemented.
    
\begin{figure}[!h]
\begin{center}
\begin{lstlisting}[language=Python]  
from nltk.corpus import wordnet as wn
...
def create_syns_from_wn(self, word):
  syns = wn.synsets(word) 
  lemmas = Set()
  for syn in syns:
    lemmas = lemmas | Set([lemma.name for lemma in syn.lemmas] )
  return lemmas

def expand_list_wn(self,  list):    
  source = Set(list)
  expanded_set = Set()
  for word in source:
    expanded_set = expanded_set | self.create_syns_from_wn(word)
return expanded_set |  source

def list_most_used_words(companies, max):
  words = flatten(map(lambda company: company.rawname.split(), companies))
  counter = collections.Counter(words)  
  return [x[0].lower() for x in 
  filter ( lambda x: x [1] > max, 
   (itertools.islice(counter.most_common(), 0, 1000)))]
\end{lstlisting}
\caption{Expanding a list of words with Wordnet syns and Counting ``most used words'' in a dataset.}
\label{figure:step-3}
\end{center}
\end{figure}

\item Dictionary-based expansion of common acronyms and filtering. A dictionary of common acronyms in corporate 
names such as ``PTY'', ``LTD'' or ``PL'' and their variants has been created by hand in order 
to be able to extract and filter acronyms. 

\item Identification of contextual information and filtering. Mainly corporate names can contain nationalities or place names that, in most of 
cases, only add noise to the real corporate name. In this case, the use of external services such as Geonames, Google Places 
or Google Maps can ease the identification of these words and their filtering. In order to tackle this situation the REST web service 
of Geonames has been selected due to its capabilities to align text with locations.

\item Spellcheking (optional). This stage seeks for providing a method for fixing mispelling errors. It is based on the 
well-known speller~\cite{NorvigSpelling} of Peter Norvig that uses a train dataset for creating a classifier. Although the accuracy of this 
algorithm is pretty good for relevant words in corporate names, the empirical and unit tests with a working dataset 
have demonstrated that mispelled non-relevant words is more efficient and accurate using a stopwords set/dictionary (this set has been 
built with words that are not in the set of ``most common words'', step 2, and exist in the Wordnet database). Furthermore some spelling corrections 
are not completely adequate for corporate names due to words could change and, therefore, a non-acceptable variant of the name 
could be accidentally included. That is why this stage is marked as optional and must be configured and performed with extreme care.

\item Pos-tagging parts of speech according to a grammar and filtering the non-relevant ones. The objective 
of this stage lies in ``classifying words into their parts of speech and labelling them accordingly is known as part-of-speech tagging''~\cite{LoperBird02}. In order 
to perform this task both a lemmatizer based on Wordnet and a grammar for corporate names (``NN''-nouns and ``JJ''-adjectives connected with articles and 
prepositions) have been designed, see Figure~\ref{figure:step-4}. Once words are tagged next step consists in filtering non-relevant categories 
in corporate names keeping nouns and adjectives, as an example Figure~\ref{figure:step-4} also shows how to walk and filter nodes in the parsed tree.

\begin{figure}[!h]
\begin{center}
\begin{lstlisting}[language=Python]  
self.lemmatizer = nltk.WordNetLemmatizer()
self.grammar = r"""
  NBAR: {<NN.*|JJ>*<NN.*>}   
  NP: {<NBAR>}
      {<NBAR><IN><NBAR>} 
    """
self.chunker = nltk.RegexpParser(self.grammar)

def leaves(self, tree):
  for subtree in tree.subtrees(filter = lambda t: t.node=='NP'):
    yield subtree.leaves()
\end{lstlisting}
\caption{Regular expression-based chunker in Python NLTK and Filtering words by the category ``NP'' (noun phrase) .}
\label{figure:step-4}
\end{center}
\end{figure}   

\item Cluster corporate names. This task is in charge of grouping names by similarity applying 
a string comparison function. Thus if the clustering is appied $n$ times any name 
will be grouped by ``the most probably/used name'' according to a threshold generated 
by the comparison function. This first version of CORFU has used the WRatio function to 
compare strings (available in the Levenshtein Python package) and a custom clustering 
implementation.

\item Validate and reconcile the generated corporate name via an existing reconcile service (optional). This last 
step has been included with the objective of linking the final corporate name with an existing information 
resource and adding new alternative labels. The OpenCorporates and DBPedia reconciliation services have 
been used in order to retrieve an URI to new corporate name. As a consequence the CORFU unifier is partially supporting 
one of the main principles of the LOD initiative such as unique identification.
\end{enumerate}
% 
%
\subsection{Use Case: Unifying supplier names in the Australian e-Procurement domain}
% % 
As previous sections have introduced there is an increasing interest and commitment 
in public bodies to create a real transparent public administration. In this sense 
public administrations are continuously releasing relevant data in different domains 
such as tourism, health or public procurement with the aim of easing the implementation 
of new added-value services and improve their efficiency and transparency. In the particular 
case of public procurement, main and large administrations have already made publicly available the information 
with regards to public procurement processes. In this case of study the information of Australia is used to test the CORFU unifier. It is comprised of a dataset 
of more than $400$K supplier names during the period 2004-2012. In order to be able to extract 
good statistics from this dataset the unification of names must be applied to. That is why 
the CORFU stepwise method has been customized to deal with the hetereogenities of this large 
dataset as Table~\ref{config-corfu} summarizes.


\begin{table}[!htb]
\renewcommand{\arraystretch}{1.3}
\begin{center}
\begin{tabular}{|p{1cm}|p{6cm}|p{6cm}|}
\hline
  \textbf{Step} & \textbf{Name} & \textbf{Customization}  \\  \hline
  $1$ & Normalize raw text and remove duplicates & Default \\ \hline
  $2$ & Filter the basic set of common stopwords in English & Default\\ \hline
  $3$ & Filter the expanded set of most common words in the dataset & Two stopwords sets: $355$ words (manually) and words with more than $n=50 $ apparitions (automatically) \\ \hline
  $4$ & Dictionary-based expansion of common acronyms and filtering & Set of $50$ acronyms variations (manually) \\ \hline
  $5$ & Identification of contextual information and filtering & Use of Geonames REST service\\ \hline
  $6$ & Spellcheking (optional) & Train dataset of $128457$ words provided by Peter Norvig's spellchecker~\cite{NorvigSpelling}. \\ \hline
  $7$ & Pos-tagging parts of speech according to a grammar and filtering the non-relevant ones & Default \\ \hline
  $8$ & Cluster corporate names & Default \\ \hline
  $9$ & Validate and reconcile the generated corporate name via an existing reconcile service (optional) & Python client and Google Refine \\ \hline
  \hline
  \end{tabular}
  \caption{Customization of the CORFU technique for Australian supplier names.}
  \label{config-corfu}
  \end{center}
\end{table} 

\section{Evaluation}
\subsection{Research design}
Since the CORFU approach has been designed and implemented it is necessary to 
establish a method to assess quantitatively the quality of the results.  The steps to carry 
out this experiment are: 1) Configure the CORFU technique, see Table~\ref{config-corfu};  
2) Execute the algorithm taking as a parameter the file containing the whole dataset of company names;
3) Validate (manually) the dump of unified names; 4) Calculate the values of 
precision and recall. In particular, this evaluation considers the precision of the algorithm as ``the number of supplier 
names that have been correctly unified under the same name'' while recall is 
``the number of supplier names that have not been correctly classified under the same name''. 

\subsection{Sample}
In order to validate the CORFU approach the dataset of supplier names in Australia 
in the period 2004-2012 containing $430188$ full names 
and $77526$ unique names has been selected. The experiment has been 
carried out executing the aforementioned steps in the whole dataset 
to finally generate a dump containing for each supplier the 
rawname and the unified name. These mappings has been validated 
by hand to quantify the typical measures of precision and recall. 

%awk -F"\t" '{print $2}' supplier-unified-names.csv | sort -n | uniq -c | sort -rnk1
%awk -F"\t" '{print $1 ";1"}' supplier-unified-names.csv > original-counters.csv

\begin{figure}[!htb]
\centering
	\includegraphics[width=6cm]{./imgs/corfu-stats}
\caption{Graphical view of supplier and number of appeareances.}
\label{fig:results-graph-3}
\end{figure}

\subsection{Results and Discussion}
According to the results presented in Table~\ref{tabla:aus-results}, the precision 
and recall of the CORFU technique are consider acceptable for the whole dataset 
due to $37248 (77526-40278)$, a $48\%$ of the supplier names, has been unified with 
a precision of $0.762$ and recall of $0.237$ (best values must be close to $1$). 

In order to improve the results for relevant companies, the experiment has also 
been performed and evaluated for the first 100 companies in the Forbes list. In this case, 
results show a better performance in terms of precision and recall and all these supplier names 
were unified by a common right name. The explanation of this result 
can be found due to some of the paramenters of the CORFU technique were 
specially selected for unifying these names because of their relevance in world economic activities.

On the other hand, it is important to emphasize that the last step of linking these names 
with existing web information resources using the reconciliation service of 
OpenCorporates or DBPedia in Google Refine has the potential of generating $28383$ right 
links ($36.61$\%) instead of the initial $8$\%. Thus the initial problem of linking $n$ string literals $\to$ 
$1$ company $\to$ $1$ URI has been substantially improved. 

Finally, the frequency distribution of supplier and number of appeareances are 
depicted on Figure~\ref{fig:results-graph-3} with the objective of presenting 
how the cloud of points (appeareances) that initially were only one per supplier has 
emerged due to the unification of names, for instance in the case of ``Oracle'' $75$ apparitions can 
now be shown.

\begin{table}[!h]
\renewcommand{\arraystretch}{1.3}
\begin{center}
\begin{tabular}{|p{3cm}|p{3cm}|p{2cm}|p{2cm}|l|l|l|}
\hline
  \textbf{Total number of companies} & \textbf{Unique names}& \textbf{CORFU unified names}& \textbf{\% of unified names} & \textbf{Precision} & \textbf{Recall} & \textbf{F1 score} \\  \hline
   $430188$ & $77526$ & $40277$  &$48\%$ & $0.762$ & $0.311$&$0.237$ \\ \hline   
   $430188$ & $X of 77526$ & $Y$ & $85\%$&  $1$ & $1$ &$1$\\ \hline
  \hline
  \end{tabular}
  \caption{Results of applying the CORFU approach to the Australian supplier names.}
  \label{tabla:aus-results}
  \end{center}
\end{table} 



\section{Conclusions and Future Work}
FIXME: The main conclusion of this works lies in the design of a technique to prepare 
raw organization names in a specific context, e.g. Australia supplier names,  
before performing a reconciliation process. Although the percentage of potential 
right links to existing datasets has been dramatically improved it is clear that 
human-validation is also required to ensure the correct unification of names. 
Other NLP techniques based on n-grams or a classifier will be used in the future 
to deliver a complete and intelligent company unifier. On the other hand, the 
application of this technique enables the comparison of rewarded contracts to 
different companies and can help to improve the transparency in public 
administrations.

\bibliographystyle{plain}
% %\bibliographystyle{unsrt}
% %\bibliographystyle{acm}
\bibliography{references}
% \renewcommand{\bibname}{References}
\end{document}

